\section[R mediation and sensitivity contracts]{R mediation and sensitivity: preserving the intervention, the mixture and the unit of uncertainty}
\label{sec:exec25-rmed}

A bank considering a servicing intervention needs to distinguish an improvement
caused by the intervention from an improvement transmitted through a particular
mechanism. Let $A$ denote an offer of restructuring, $M$ subsequent payment
engagement, and $Y$ loss over a fixed horizon. Baseline information $W$ precedes
the offer. A later liquidity assessment $Z$ may be changed by the offer and
affect both engagement and loss. The question therefore requires a declared
intervention on engagement, a target population, an information clock and
assumptions linking the intervention to observed data. R packages supply useful
approaches, but their targets and uncertainty calculations are not interchangeable.
This comparison extends the preceding mediation, medflex, sensemakr and EValue
discussions with CMAverse, medoutcon, regmedint, causalsens, tipr and BayesGmed.
Inspected sources and authored examples support conditional reuse decisions;
they do not establish native qualification. \citep{exec25_rmed_choice}

\subsection{A causal target before a software method}

Write $Y_{a,m}$ for the outcome after setting offer and mediator to $(a,m)$,
and $M_a$ for the mediator under offer $a$. Two complementary natural
decompositions on an additive scale are
\[
 \operatorname{TE}=E(Y_{1,M_1}-Y_{0,M_0})
 =\operatorname{PNDE}+\operatorname{TNIE}
 =\operatorname{TNDE}+\operatorname{PNIE}.
\]
Here $\operatorname{PNDE}=E(Y_{1,M_0}-Y_{0,M_0})$ and
$\operatorname{TNIE}=E(Y_{1,M_1}-Y_{1,M_0})$; the other pair interchanges the
reference treatment. The controlled direct effect $E(Y_{1,m}-Y_{0,m})$ instead
holds engagement at the same specified level. It need not equal a natural
effect, and subtracting it from TE need not yield a natural indirect effect.
The selected BayesGmed primary text develops these distinctions under consistency,
positivity, no unmeasured confounding and a cross-world condition. Randomizing
the offer alone does not supply every mediation assumption; a posterior interval
does not establish them either. \citep{exec25_rmed_yimer_paper}

When offer-induced liquidity confounds engagement and loss, an alternative
draws a new mediator $G_b$ from the distribution of $M_b$ conditional on $W$,
independently of the realized intermediate confounder under the declared
intervention law. Let $\mu(a,z,m,w)=E(Y\mid A=a,Z=z,M=m,W=w)$. The selected
D\'iaz and colleagues paper identifies the corresponding mean as
\[
 \theta(a,b)=\int \mu(a,z,m,w)\,p(z\mid a,w)\,
                  p(m\mid b,w)\,p(w)\,d(z,m,w).
\]
The identifying requirements include the relevant treatment and mediator
exchangeabilities, consistency and conditional support. This intervention
differs from drawing a mediator conditional on the intermediate confounder as
well as baseline information. Changing the conditioning set changes the
pathways assigned to direct and indirect components.
\citep{exec25_rmed_diaz_paper}

Define $\operatorname{IDE}=\theta(1,0)-\theta(0,0)$ and
$\operatorname{IIE}=\theta(1,1)-\theta(1,0)$. Their sum is the total contrast
under the randomized mediator intervention, which need not equal the ordinary
treatment effect. An authored example takes
$P(Z=1\mid A=0)=1/4$, $P(Z=1\mid A=1)=1/2$,
$P(M=1\mid Z)=1/4+Z/2$, and $Y=ZM$. The ordinary treatment effect is $3/16$,
whereas IDE is $3/32$ and IIE is $1/16$, summing to $5/32$. All mediator
probabilities are positive. The difference arises from changing dependence in
the intervention, not estimation error. The product must label the two totals
explicitly. \citep{exec25_rmed_choice}

\subsection{CMAverse: a broad interface with a consequential draw contract}

CMAverse version 0.1.0 documents regression, weighting, inverse odds weighting,
natural-effect, marginal structural and g-formula approaches. Its interface
restricts methods when post-treatment confounders are supplied and changes
natural-effect labels to randomized analogues. Parameter-function and imputation
estimators have different applicability; the documented logistic and Cox
parameter-function cases require a rare outcome. Its sensitivity interface
distinguishes unmeasured confounding from measurement error and documents a
single-variable measurement-error scope. These are useful capabilities to
compare, provided the proposed product exposes method, scale and restrictions.
\citep{exec25_rmed_cmaverse}

The selected g-formula delegate simulates intermediate confounders from
exposure and baseline covariates, then mediators using those variables and the
simulated intermediate values. In its post-treatment-confounder branch, it
subsequently permutes each mediator column across eligible rows. Outcome
prediction retains the row's baseline and intermediate variables. Global
permutation requires a different justification from a draw conditional on that
row's baseline information; separate column permutations also require
justification when a joint mediator law is intended. \citep{exec25_rmed_cmaverse}

For $W\sim\operatorname{Bernoulli}(1/2)$, mediator probabilities $1/4$ and
$3/4$ in the two baseline groups, and outcome regression $WM$, the
baseline-conditioned stochastic mean is $3/8$. A global shuffle pairs each
row with the marginal mediator probability $1/2$, yielding $1/4$ in expectation.
For two identical balanced binary mediators, $E(M_1M_2)=1/2$; independently
permuting columns yields $1/4$ in expectation. These exact witnesses specify
invariants for an adapter. They are neither native failure frequencies nor a
rejection of g-computation. The route is to declare the conditional joint draw,
reconcile the delegate, and retain an explicitly different target or a separately
reviewed repair. \citep{exec25_rmed_choice}

Conditional R invocation is preferable to porting a large method collection
before resolving its target. Requests should carry mediator order, conditioning
set, treatment contrast, scale, simulation budget and random state. Resampling
units must match the data law: a row bootstrap is not automatically branch-cluster
or shared-calendar inference. Imputation, simulation integration and bootstrap
variation require separate accounting. The GPL declaration and full dependency
inventory also accompany contemplated redistribution; this source review does
not clear every artifact. \citep{exec25_rmed_cmaverse,exec25_rmed_choice}

\subsection{medoutcon: efficient scores and the sampling contract}

medoutcon targets natural and interventional effects using one-step and targeted
minimum loss estimators. The primary paper assumes independent, identically
distributed observations, derives an efficient influence function, and requires
positivity and products of nuisance-estimation errors vanishing faster than
$n^{-1/2}$. Cross-fitting relaxes empirical-process restrictions; it does not
prove those rates for arbitrary learners or remove hidden confounding and
dependence. A low-dimensional intermediate-confounder representation avoids
directly estimating a high-dimensional mediator density, a useful advantage
for multivariate engagement measurements within the stated conditions.
\citep{exec25_rmed_diaz_paper}

The selected source switches to natural effects when $Z$ is explicitly absent,
constructs binary contrasts, and rescales outcomes by their observed range.
A constant outcome makes that scaling denominator zero. Contrasts subtract
influence scores before variance calculation; propensity clipping must be
recorded as a numerical choice rather than proof of population support. For
two-phase mediator measurement, the source uses
\[
 D_{\mathrm{obs}}=\frac{R}{\pi}D
                 -\left(\frac{R}{\pi}-1\right)m_D,
\]
where $R$ indicates observation, $\pi$ is inclusion probability and $m_D$
predicts the complete-data influence quantity. With $D=3$, $m_D=1$ and
$\pi=1/2$, the observed and missing cases give 5 and 1, averaging to 3 under
that law. False inclusion probabilities do not become valid through calculation.
\citep{exec25_rmed_medoutcon,exec25_rmed_choice}

In the inspected one-step survey-weight branch the estimate is a weighted
mean, but the variance expression multiplies the influence quantity by supplied
weights without corresponding normalization. For scores $(1,3)$, changing
weights from $(1,1)$ to $(2,2)$ retains estimate 2 while changing the displayed
variance-of-the-mean expression from 1 to 4. Requiring normalized weights may
reconcile that convention; an adapter accepting arbitrary relative weights
cannot silently assume it. The proportion-mediated option is a log-ratio
functional, distinct from an additive proportion. These observations motivate
qualification tests rather than establish installed-estimator coverage.
\citep{exec25_rmed_medoutcon,exec25_rmed_choice}

The conditional adoption route retains the R learner stack until independent
fixtures establish target, score, contrast, rescaling and uncertainty parity.
Public trials should include weak support, constants, weight rescaling,
incomplete mediator measurement and deliberately invalid dependent sampling.
A port needs a deployment benefit sufficient to maintain these scientific
invariants. \citep{exec25_rmed_choice}

\subsection{regmedint: interactions and conditional effects}

regmedint supplies a more explicit parametric alternative, carrying baseline
evaluation values, mediator and outcome model families, treatment--mediator
interaction and effect modification. Its case-control branch fits the mediator
in controls and states a rare-outcome requirement. The selected delta
calculation constructs separate regression covariance blocks. Evaluation at
mean covariates is not automatically population standardization, and a block
covariance is not automatically appropriate for clustered servicing data.
\citep{exec25_rmed_regmedint}

Consider $M=1+2A+\varepsilon_M$ and
$E(Y\mid A,M)=3A+4M+5AM$, with mean-zero disturbance and the necessary causal
assumptions. PNDE and TNIE are 8 and 18; TNDE and PNIE are 18 and 8; TE is
26. The controlled direct effect at zero is 3, and the coefficient product
$2\times4=8$ gives only one indirect component. Output must retain interaction,
reference treatment and mediator setting. Proportion mediated also needs a
nonzero interpretable total and need not lie between zero and one.
\citep{exec25_rmed_regmedint,exec25_rmed_choice}

Use this package conditionally as a parametric comparator and fixture source.
When a marginal effect is needed, define the standardization explicitly.
Logistic formulas, log effects and rare-outcome approximations should be checked
against exact binary enumeration at both small and substantial baseline risks.
A port should preserve the chosen approximation rather than present an exact
general risk decomposition. \citep{exec25_rmed_choice}

\subsection{Sensitivity measures have units, directions and domains}

causalsens asks how a linear outcome analysis changes under a specified
selection-bias function. It adjusts outcomes and refits the linear model using
model-based uncertainty. Its one-sided adjustment is $\alpha(A-e)$ for
propensity $e$, whereas alignment uses
$\alpha\{(1-e)A+e(1-A)\}$. With $e=1/2$ and $\alpha=4$, the first reduces a
treated-minus-control mean gap by 4; the second subtracts 2 from both groups
and leaves the gap unchanged. Treated-population functions differ again.
Thus $\alpha$ carries outcome units and a confounding direction. A sensitivity
curve does not identify a mediation mechanism. \citep{exec25_rmed_causalsens,exec25_rmed_choice}

tipr provides coefficient adjustment and tipping-point interfaces. Its selected
partial-$R^2$ and robustness functions delegate to sensemakr. A simple coefficient
adjustment subtracts a declared bias product. Its numeric E-value helper applies
$r+\sqrt{r(r-1)}$, after reciprocating an input at or below one, without a
baseline-risk or effect-scale argument. An odds ratio of 2 at baseline risk
$1/2$ corresponds to risk ratio $4/3$. Applying the formula to 2 gives about
3.414, whereas applying it to $4/3$ gives 2. This exact example rejects silent
scale substitution; it does not supply a universal hazard-ratio conversion.
Nonpositive inputs and intervals crossing the null need explicit handling.
\citep{exec25_rmed_tipr,exec25_rmed_choice}

The inspected boundary helper contains a fallback that resets outcome
partial-$R^2$ when an interval misses its zero-boundary check. An adapter must
preserve the requested scenario and verify its endpoint, reporting a failed
boundary calculation when appropriate. Conditional interface reuse can build
on the previously reviewed underlying methods while keeping their scientific
and numerical contracts visible. \citep{exec25_rmed_tipr,exec25_rmed_choice}

\subsection{BayesGmed: verify the latent-state likelihood before sampling}

BayesGmed offers Bayesian parametric g-computation for continuous or binary
mediators and outcomes. Its selected primary text describes independent
observations, identifying assumptions, correctly specified conditional models,
posterior simulation and bootstrap integration over baseline covariates.
The sensitivity discussion introduces a binary unobserved confounder with a
prevalence model and bias parameters. The paper describes uniform bias priors;
selected version 0.0.3 instead supplies a multivariate Gaussian prior through a
covariance matrix. These are different analyses requiring explicit justification.
The wrapper's link arguments also do not alter its fixed binary/continuous
dispatch. Unsupported link requests should be rejected rather than reported
as model changes. \citep{exec25_rmed_yimer_paper,exec25_rmed_bayesgmed}

More fundamentally, if $p=P(U=1\mid A)$ and $f_u(m,y)$ is the conditional
mediator--outcome density in state $u$, the normalized mixture is
\[
 L(m,y)=(1-p)f_0(m,y)+p f_1(m,y).
\]
All four inspected Stan sensitivity files attach $\log(1-p)$ to both latent
states before log-sum-exp. The two inspected generated C++ counterparts retain
the same expression. It therefore becomes
\[
 L_{\mathrm{src}}(m,y)=(1-p)\{f_0(m,y)+f_1(m,y)\},
 \qquad \int L_{\mathrm{src}}=2(1-p).
\]
This is a static normalization discrepancy at the selected revision, not a
measured sampler failure rate. It concerns the inspected sensitivity paths
and does not by itself reject ordinary models or Bayesian mediation generally.
\citep{exec25_rmed_bayesgmed}

An exact binary witness takes $p=4/5$, mediator probabilities $(1/4,3/4)$ and
outcome probabilities $(1/5,4/5)$ in states zero and one, conditionally
independent within each state. The intended probability of $(M,Y)=(1,1)$ is
$49/100$. The source expression gives $13/100$ and sums to $2/5$ over the
four observations. Renormalization gives $13/40$, still different because it
substitutes equal mixing weights. When the components coincide, the intended
likelihood loses dependence on $p$, but the source expression retains it.
These checks should precede expensive sampling. \citep{exec25_rmed_choice}

Defer the selected sensitivity implementation until mixture, prior, link and
artifact semantics are reconciled. A repair needs a versioned successor,
independent likelihood enumeration, posterior and predictive checks,
integration diagnostics and public known-law comparisons. Ordinary-model
review can proceed independently. Rights also need clarification: package
metadata declares MIT while a retained separate Stan license include states
GPL3 or later. Its applicability to a redistributed artifact is unresolved;
neither an installed binary nor all dependency rights is qualified here.
\citep{exec25_rmed_bayesgmed,exec25_rmed_choice}

\subsection{A common contract and public qualification route}

A CLI or MCP request should name baseline and post-treatment variables,
mediator ordering and joint law, intervention conditioning, treatment values,
target population, outcome scale and horizon, sampling units, measurement
inclusion law and sensitivity parameters with their units. Responses should
distinguish unsupported identification, support failure, undefined ratios,
numerical failure and an admissible estimate. They should retain exact source
revision, method, prior, folds, resampling unit, integration budget and
uncertainty definition. An LLM can help elicit declarations and present
alternatives, but cannot infer scientific truth or choose a conditioning set
merely to make a backend run. \citep{exec25_rmed_choice}

Public qualification starts with finite known-law examples, followed by frozen
simulations with supported and intentionally unsupported structures. Parametric
formulas, semiparametric scores, sensitivity wrappers and Bayesian likelihoods
can be tested independently. Complete request counts, refusals, target errors,
interval behavior and end-to-end costs are needed. For hypothetical USD 30,000
spent on 1,000 requests, 600 returned results and 400 scientifically admissible
results, costs are USD 30, USD 50 and USD 75 per respective denominator. Runtime
alone cannot select a denominator or establish bank value. \citep{exec25_rmed_choice}

Four bounded repository queries returned 105 unique leads, with unreturned
and uninspected tails. They do not establish an exhaustive R census. Source
inspection, artifact and rights resolution, and public native experiments
remain necessary. Unknown cross-world relationships, hidden mediator
confounding and unrestricted dependence need information or justified
restrictions. These distinctions preserve a practical route to R mediation
and sensitivity qualification. \citep{exec25_rmed_choice}
